% !PRINTED 635
\weberpage{635}
Da wir jetzt im Stande sind, für jedes $s$ die Primfactoren von $\psi_s(r)$ zu ermitteln, gehen wir dazu über, die Zahl $(r,\eta)^m$ nach der Formel (11) in ihre Primfactoren zu zerlegen.

Wir bilden zu diesem Zwecke nach (15) die kleinsten positiven Reste $\varkappa$ der Zahlen, $-n, -2n, \cdots -(m-2)n$ nach dem Modul $m$, und finden, von der Ordnung abgesehen, die Zahlen
\[
\varkappa = 1, 2, \ldots, n-1, n+1, \ldots, m-1;
\]
die Reste $0$ und $n$ kommen darunter nicht vor, weil $s$ keinen Werth erhält, für den $s$ oder $s+1$ durch $m$ theilbar wird. Von diesen Zahlen sind aber gerade $n-1$ kleiner als $n$, und so oft kommt also der Factor $\mathfrak{p}_{n'}$ in dem Producte $\psi_1\psi_2\ldots\psi_{m-2}$ vor. Dann enthält $p$ diesen Factor noch einmal und folglich kommt er nach (11) genau $n$ mal in $(r,\eta)^m$ vor. Da keine anderen als die Primfactoren $\mathfrak{p}_{n'}$ in $(r,\eta)^m$ aufgehen können, so ist hierdurch die Zerlegung vollständig ausgeführt:
\weq{16}{(r,\eta)^m = \prod^{n} \mathfrak{p}_{n'}^{\,n};}
hierin durchläuft $n$ die Reihe der positiven Zahlen, die kleiner als $m$ und relativ prim zu $m$ sind, und $n'$ ist jedesmal aus der Congruenz
\weq{17}{n n' \equiv 1\ (\mathrm{mod}\ m)}
zu bestimmen\footnote{Dieser Satz rührt von \gesperrt{Kummer} her. Vgl. Theorie der idealen Primfactoren der complexen Zahlen etc. Abhandlungen der Berliner Akademie, 1856.}.

Die Formel (16) gilt für ein gerades wie für ein ungerades $m$ und ist auch noch für den Fall $m=2$ gültig, für den sie ein schon bekanntes Resultat giebt, da dann $(r,\eta) = (-1,\eta)$ mit einer \gesperrt{Gauss}'schen Summe übereinstimmt (Bd.\ I, \S.\,171).

Wir wollen aber nun, indem wir den Fall $m=2$ jetzt ausschliessen, der Formel (16) eine etwas allgemeinere Gestalt geben.

Es sei $m=q^{\varkappa}$ eine beliebige Primzahlpotenz, nur grösser als $2$, und $p$ eine von $q$ verschiedene Primzahl von der Eigenschaft, dass $p-1$ durch $q$ oder, wenn $q=2$ ist, durch $4$ theilbar ist.

Es sei ferner $m_1$ der grösste gemeinschaftliche Theiler von $p-1$ und $m$, und
\[
m = m_1 m_2,
\]
%CONT>>
